\subsection{Transfer of component error into retrieval error}\label{sec:transfer}

Section~\ref{sec:stability} controls the inverse where the denominator stays away from zero. The theorem below
describes a whole evaluation distribution, with entries of vanishing transmission included, with the constant one,
and shows that the resulting loss of rate is not an artefact of the argument. Let $P$ be the evaluation distribution
over entries (a state, a band, and the surface reflectance when it is varied), in fixed physical units; expectations
are conditional on the fitted emulator. Assume
\begin{equation}\label{eq:transfer-hyp}
0\le\rho\le R\le1,\qquad t>0,\qquad 0\le s\le S,\qquad RS<1 .
\end{equation}
The emulator predicts $(\widehat a,\widehat t,\widehat s)$. The \emph{constrained retrieval} replaces $\widehat t$ by
$\max\{\widehat t,0\}$ and $\widehat s$ by its projection onto $[0,S]$, neither of which increases $|e_t|$ or
$|e_s|$, and sets
\begin{equation}\label{eq:rhobar}
\bar\rho=\Pi_{[0,R]}\Bigl(\frac{L-\widehat a}{\widehat t+\widehat s\,(L-\widehat a)}\Bigr)
\end{equation}
when the denominator is positive, and otherwise returns a prescribed value in $[0,R]$ and records a failure. It is
an estimator in its own right and is reported beside the unconstrained inverse and its failure counts, never in
place of them. The \emph{retrieval-weighted component error}, computed after the projections, is
\begin{equation}\label{eq:eR}
e_R=|e_a|+R\,|e_t|+R^2\,t\,|e_s|,\qquad \varepsilon_R^2=\E_P[e_R^2],
\end{equation}
which is at most $e=|e_a|+|e_t|+t|e_s|$ since $R\le1$; $\varepsilon_R$ is not the mean relative radiance error of
the tables. Put $F_t(\tau)=\Prob(t\le\tau)$.

\begin{theorem}[error transfer with constant one]\label{thm:transfer}
Under \eqref{eq:transfer-hyp}, for the constrained retrieval \eqref{eq:rhobar}:
\begin{enumerate}
\item $|\bar\rho-\rho|\le\min\{R,\;e_R/t\}\le\min\{R,\;e/t\}$ at every entry, and no universal constant below one
is possible.
\item For every $\tau>0$,
\[
\E_P|\bar\rho-\rho|^2\le\E_P\min\Bigl\{R^2,\frac{e_R^2}{t^2}\Bigr\}
\le R^2F_t(\tau)+\E_P\Bigl[\frac{e_R^2}{t^2}\1\{t>\tau\}\Bigr]
\le R^2F_t(\tau)+\frac{\varepsilon_R^2}{\tau^2}.
\]
\item If $F_t(u)\le Au^\beta$ for $0<u\le u_0$ with some $\beta>0$, then for
$\varepsilon_R\le u_0^{(\beta+2)/2}$,
$\E_P|\bar\rho-\rho|^2\le(R^2A+1)\,\varepsilon_R^{2\beta/(\beta+2)}$; a component-error bound
$\varepsilon_{R,n}=O(n^{-r})$ yields $\|\bar\rho_n-\rho\|_{L^2(P)}=O(n^{-r\beta/(\beta+2)})$.
\item If $t\ge t_0>0$ $P$-almost surely, then $\|\bar\rho-\rho\|_{L^2(P)}\le\varepsilon_R/t_0$.
\end{enumerate}
\end{theorem}

\noindent\emph{Proof in Section~\ref{app:proofs}.}
The middle expression in (ii) keeps the association between the errors and the transmission and is the one the
tables evaluate; the last needs no power law. Both are computable on a finite test table; (iii) is a population
statement whose tail hypothesis a finite table cannot establish down to $t=0$. The exponent in (iii) cannot be
improved from its hypotheses. Item (i) holds entry by entry, so it passes to quantiles: on any set of entries, a
single band for instance, every quantile of $|\bar\rho-\rho|$ is at most the same quantile of $\min\{R,e_R/t\}$,
which depends only on the component errors and the true transmission. In Bayesian inversion the corresponding
statements bound the Hellinger distance between the true posterior and one built on a Gaussian-process emulator by
moments of the emulator's error \cite{stuartteckentrup}; here the inverse is deterministic and the weight of each
entry, $1/t$, is explicit.

\begin{proposition}[the exponent is attained]\label{prop:attained}
For every $\beta>0$, every $0<c<r_0<R$ and every $\tau\in(0,1)$ there are an evaluation distribution and a predictor
satisfying \eqref{eq:transfer-hyp} with $S=0$ and $F_t(u)=u^\beta$ on $(0,1)$, namely $\rho=r_0$, $a=s=0$,
$\widehat t=t$, $\widehat s=0$ and $\widehat a=ct\,\1\{t\le\tau\}$, for which $\bar\rho=r_0-c\,\1\{t\le\tau\}$,
\[
\begin{aligned}
\varepsilon_R^2&=\frac{c^2\beta}{\beta+2}\,\tau^{\beta+2},
&\E_P|\bar\rho-\rho|^2&=c^2\tau^\beta,\\
\E_P|\bar\rho-\rho|^2
&=c^{4/(\beta+2)}\Bigl(\frac{\beta+2}{\beta}\Bigr)^{\beta/(\beta+2)}
\varepsilon_R^{2\beta/(\beta+2)} .
\end{aligned}
\]
\end{proposition}

\noindent\emph{Proof in Section~\ref{app:proofs}.}
This is sharpness of the transfer inequality for the physical inverse, not a minimax statement over retrieval
algorithms. What decides the rate is where the errors sit relative to the transmission.

\begin{theorem}[conditional error profile]\label{thm:profile}
Under \eqref{eq:transfer-hyp}, suppose $\E_P[e_R^2\mid t]\le\sigma^2t^{2\nu}$ $P$-almost surely for some $0\le\nu\le1$. Then
\[
\E_P|\bar\rho-\rho|^2\le\E_P\min\{R^2,\;\sigma^2t^{-2(1-\nu)}\} .
\]
Under the tail hypothesis of Theorem~\ref{thm:transfer}(iii) and $\nu<1$, as $\sigma\to0$,
\[
\E_P|\bar\rho-\rho|^2=
\begin{cases}
O(\sigma^{\beta/(1-\nu)}), & \beta<2(1-\nu),\\
O(\sigma^2\log(1/\sigma)), & \beta=2(1-\nu),\\
O(\sigma^2), & \beta>2(1-\nu),
\end{cases}
\]
and at $\nu=1$, $\E_P|\bar\rho-\rho|^2\le\sigma^2$ with no tail hypothesis. In the model of
Proposition~\ref{prop:attained} the predictor $\widehat a=a+\min\{ct,\sigma t^\nu\}$ has $e_R\le\sigma t^\nu$ and
$|\bar\rho-\rho|=\min\{c,\sigma t^{\nu-1}\}$, and attains each order.
\end{theorem}

\noindent\emph{Proof in Section~\ref{app:proofs}.}
At $\nu=0$ the theorem recovers the uniform-error regimes under the weaker hypothesis of a bounded conditional
second moment; at $\nu=1$ an emulator whose conditional root-mean-square error decreases proportionally to the
transmission loses no rate. The central diagnostic is therefore the profile $t\mapsto\E[e_R^2\mid t]$ and not only
the tail $t\mapsto F_t(t)$: the same overall error concentrated in the opaque entries gives the worst-case loss,
and proportional to the transmission gives none.

\subsection{The stacking objective}\label{sec:joint}

\paragraph{Affine fitting predictions and projected retrieval.}
Let $w$ denote simplex weights, with one simplex for each component being combined.
Write $\widetilde a_w$, $\widetilde t_w$ and $\widetilde s_w$ for the affine combinations
used in fitting, and $\widetilde e_a=\widetilde a_w-a$, $\widetilde e_t=\widetilde t_w-t$,
$\widetilde e_s=\widetilde s_w-s$ for their errors. Retrieval uses
$\widehat t_w=\max\{\widetilde t_w,0\}$ and
$\widehat s_w=\Pi_{[0,S]}\widetilde s_w$, followed by \eqref{eq:rhobar}.
Base predictions may instead be projected before combination, provided the values
entering the fitting objective remain affine in $w$. In either case, projection during
retrieval does not increase flux or albedo error under \eqref{eq:transfer-hyp}.

For a fixed finite threshold $\tau>0$, the total-flux raw-error surrogate is
\begin{equation}\label{eq:Jtau}
J_\tau(w)=\E_P\!\left[
\frac{\widetilde e_a^2+R^2\widetilde e_t^2+R^4t^2\widetilde e_s^2}
     {(t\vee\tau)^2}\right].
\end{equation}
It is a convex quadratic: the errors are affine in the weights and the entry weights
are nonnegative and independent of $w$. Its empirical version replaces the expectation
by the validation average. Theorem~\ref{thm:transfer} and Cauchy--Schwarz imply
\begin{equation}\label{eq:Jtau-bound}
\E_P|\bar\rho_w-\rho|^2\le R^2F_t(\tau)+3J_\tau(w).
\end{equation}
Indeed, projection gives
$e_R\le|\widetilde e_a|+R|\widetilde e_t|+R^2t|\widetilde e_s|$; split the
expectation at $t=\tau$ and apply the three-term square bound on $\{t>\tau\}$.
This inequality holds for every feasible $w$, not only an optimizer.

\paragraph{A separable surrogate.}
The stacks of Table~\ref{tab:v2-stacks} combine $Y_2$ and $Y_3$ with separate weight
vectors. Each base albedo is projected onto $[0,S]$ before fitting, and the base fluxes are
not projected. With
$\widetilde e_j=\widetilde Y_j-Y_j$ for $j=2,3$, their finite-threshold fitting objective is
\begin{equation}\label{eq:Jtau-sep}
J'_\tau(w)=\E_P\!\left[
\frac{\widetilde e_a^2+R^2\widetilde e_2^2+R^2\widetilde e_3^2
      +R^4t^2\widetilde e_s^2}{(t\vee\tau)^2}\right].
\end{equation}
This is also a convex quadratic. It omits the cross term
$2R^2\widetilde e_2\widetilde e_3$ present in $J_\tau$.
The flux floor is applied after fitting, to the combined total flux when the retrieval is
scored, and is not part of the fitted loss.
Since $t\ge0$,
\[
|\max\{\widetilde Y_2+\widetilde Y_3,0\}-t|
 \le|\widetilde e_2+\widetilde e_3|
 \le|\widetilde e_2|+|\widetilde e_3|.
\]
The same argument now uses four squared terms and gives
\begin{equation}\label{eq:Jtau-sep-bound}
\E_P|\bar\rho_w-\rho|^2\le R^2F_t(\tau)+4J'_\tau(w).
\end{equation}
For the separable fit the factor is therefore four rather than three.
Both statements require the physical hypotheses on the distribution being averaged;
full-table statistics that include entries outside those hypotheses are not certificates
of either bound.

Putting projection \emph{inside} a squared fitting loss is a different problem and need
not preserve convexity. On $0\le w\le1$, for example,
$g(w)=(\max\{2w-1,0\}-1)^2$ satisfies
$g(1/2)=1>(g(0)+g(1))/2=1/2$. The raw-error quadratic \eqref{eq:Jtau-sep} is not of
this kind.

\paragraph{Forward-optimal weights.}
A stack fitted on a forward objective weights every band alike, while the retrieval divides the
error at each entry by its transmission. The following example shows how far apart the two
optima can be.

\begin{proposition}[forward-optimal and retrieval-optimal weights]\label{prop:stackweights}
Let $0<\rho<R\le1$. For a convex combination of heads with weights $w$, let $F(w)$ be the mean
squared radiance error at $\rho$ and $E(w)$ the mean squared error of the constrained retrieval
\eqref{eq:rhobar}, and let $w_f$ minimize $F$ and $w_r$ minimize $E$.
\begin{enumerate}
\item Take two equally weighted entries with $a=s=0$ and transmissions $t_1=1$,
$t_2=\delta\in(0,1]$, and two heads that predict $t$ and $s$ exactly and the path radiance with
errors $e^A=(0,c\delta)$ and $e^B=(b,0)$, where $0<b,c\le\min\{\rho,R-\rho\}$; $w$ is the weight on
$A$. Then
\[
F(w)=\tfrac12\bigl[(1-w)^2b^2+w^2c^2\delta^2\bigr],\qquad
E(w)=\tfrac12\bigl[(1-w)^2b^2+w^2c^2\bigr],
\]
$w_f=b^2/(b^2+c^2\delta^2)$, $w_r=b^2/(b^2+c^2)$, $F(w_f)\le F(w_r)$, and
\[
\frac{E(w_f)}{E(w_r)}=\frac{(b^2+c^2\delta^4)(b^2+c^2)}{(b^2+c^2\delta^2)^2},
\]
which tends to $1+c^2/b^2$ as $\delta\to0$ and equals $(1+\delta^2)^2/(4\delta^2)$ at
$c=b/\delta$. For every $\tau\le\delta$ the transmission-weighted objective $J_\tau$ of
\eqref{eq:Jtau} equals $E$.
\item Suppose instead that $s=0$ at every entry, that every head predicts $t$ and $s$ exactly,
and that the retrieval of every combination lies in $[0,R]$. Let $\mathcal L$ be the linear span of
the heads' path-radiance errors, and let $\lambda_{\max}$ and $\lambda_{\min}$ be the supremum and the
infimum of $\E_P[r^2/t^2]/\E_P[r^2]$ over $r\in\mathcal L$ with $\E_Pr^2>0$. Then
\[
E(w_f)\le\frac{\lambda_{\max}}{\lambda_{\min}}\,E(w_r),\qquad
\frac{\lambda_{\max}}{\lambda_{\min}}\le\Bigl(\frac{t_{\max}}{t_{\min}}\Bigr)^2
\ \text{ if }\ t_{\min}\le t\le t_{\max}\ P\text{-almost surely.}
\]
\end{enumerate}
\end{proposition}

\noindent\emph{Proof in Section~\ref{app:proofs}.}
The weights that minimize the forward error put almost all their mass on the head whose error
sits at the opaque entry, because there it costs little radiance; the retrieval divides that
error by $\delta$ and pays for it in full. With $c/b$ large the forward-optimal stack is better
than either head in radiance and worse than the retrieval-optimal stack by any factor in the
inverse, which is the pattern of Table~\ref{tab:tails} in miniature. Item (ii) says when this can
happen. The ratio $\lambda_{\max}/\lambda_{\min}$ measures how unequally the transmission weights
the directions of error available to the combination: it equals one when every error in
$\mathcal L$ has the same transmission profile, and it is large only when the heads put their
errors at different transmissions, as the two heads of item (i) do. In that example it equals
$\delta^{-2}$, and the ratio $E(w_f)/E(w_r)$ at $c=b/\delta$ comes within a factor of four of it.
Without such structure the bound falls back on the squared dynamic range of the transmission.
On the physical domain of the EMIT table the transmission itself spans about twenty-four decades,
and thirteen above the floor $10^{-12}T_{\rm train}$ of Section~\ref{sec:tq-domain}, so the
factor is of order $10^{48}$ and $10^{26}$. The bound \eqref{eq:Jtau-bound}
for $J_\tau$ involves neither quantity.

\paragraph{Numerical weights and threshold convention.}
The weights are computed by nonnegative least squares with a heavily weighted row penalizing
$\sum_h w_h-1$, and are then normalized to sum to one. This gives feasible
simplex weights, though not an exact solution of the equality-constrained quadratic program,
and the mean-relative-norm arm runs six iteratively reweighted least-squares sweeps from the
squared-loss solution. The weights reported below are therefore approximate minimizers.

A finite threshold is written $\tau=uT_{\rm train}$, where $T_{\rm train}>0$ is a flux scale of
the training data, and Table~\ref{tab:v2-stacks} reports $u$. The unweighted fit is a plain
squared-error baseline on the physical components, not the limit of \eqref{eq:Jtau-sep} as
$\tau\to\infty$. The ladder of thresholds is a sensitivity analysis; a deployed threshold would
be fixed in advance by a coverage rule. The stack weights here are fitted on validation rows
that also serve the other model decisions, so Proposition~\ref{prop:selection} does not apply to
them.
